\section{Introduction}
Inferring coupling strength between dynamical systems is difficult
when shared forcing synchronizes their dynamics. Once two oscillators
move in lockstep, observations from the coupled system look the same
as observations from two independent systems driven by the same
external signal---the coupling becomes invisible. This confound arises
in neural systems exhibiting coherent oscillations~\citep{pikovsky2001},
ecological communities driven by shared environmental
cycles~\citep{grenfell1998}, and spatially connected epidemics
synchronized by seasonal forcing. Here we study this problem in the
context of infectious disease epidemiology, where abundant
surveillance data and well-characterized transmission dynamics make it
possible to quantify exactly when coupling can and cannot be inferred.

The natural approach---fitting a mechanistic transmission model to
case data from two or more regions and estimating connectivity as a
model parameter---has a long history in epidemiology and related
fields.  Time-series methods such as Convergent Cross
Mapping~\citep{sugihara2012} have been applied to identify
environmental drivers of influenza~\citep{deyle2016} and climatic
modulation of dengue~\citep{rypdal2019}. Likelihood-based approaches
have been used to estimate pathogen interaction
strengths~\citep{shrestha2011} and spatial coupling
parameters~\citep{kiss2024towards, kiss2023parameter}. Outside
epidemiology, analogous strategies are routine: neuroscientists infer
functional connectivity from synchronized neural
recordings~\citep{reid2019, friston2011}, ecologists estimate dispersal from
spatially replicated population time
series~\citep{bjornstad1999spatial}, and climate scientists quantify
teleconnections from correlated regional
indices~\citep{tsonis2004architecture}.

Yet results from these approaches have been mixed, and the reasons are
instructive. \citet{cobey2016} found that CCM was ``extremely
sensitive to periodic fluctuations, inferring interactions between
independent strains that oscillate with similar frequencies.'' The
developers of CCM acknowledged that such findings may reflect
``mistaking synchrony for causation''~\citep{sugihara2017}.
\citet{shrestha2011} found that some pathogen-interaction parameters
were ``extremely well identified'' while others remained poorly
constrained despite extensive data. In spatial epidemiology,
\citet{welch2011} showed that transmission rate and network structure
are fundamentally confounded, making it ``difficult to distinguish
between a mildly transmissible disease on a highly connected network
from a highly transmissible infection on a sparse network.'' These
results point to a common failure mode: when systems are synchronized,
an analyst can obtain a point estimate of coupling that looks
perfectly reasonable yet carries no statistical support.  The
optimizer converges, the fit looks good, but the confidence interval
spans the entire parameter space.

Most spatial epidemic studies have sidestepped this problem by relying
on external mobility data---commuting surveys~\citep{viboud2006},
gravity models~\citep{charu2017}, or mobile phone
records~\citep{kraemer2020}---rather than estimating connectivity from
case counts. Yet such proxy measures may not accurately reflect
disease transmission pathways, and are not universally
available~\citep{wu2013sensitivity}. A principled understanding of
\emph{when} connectivity can be inferred from surveillance data---and
when it cannot---would help researchers avoid wasted effort and,
more importantly, prevent publication of spurious estimates.

The Cram\'er-Rao Lower Bound (CRLB) provides exactly this kind of
guarantee. Derived from the Fisher Information
Matrix~\citep{casella2024statistical}, the CRLB gives a lower bound
on the variance of any unbiased estimator~\citep{wilke2020}. If the
CRLB for a parameter is large relative to the parameter itself, then
no estimation method---however sophisticated---can produce a reliable
estimate from the data at hand. Structural identifiability analysis
can determine whether parameters are in principle recoverable from
ideal data~\citep{tuncer2018, chowell2023structural}, but it provides
only a binary answer. The CRLB quantifies \emph{how much} information
the data actually contain, giving a practical assessment that accounts
for sample size, noise level, and the specific parameter values.
Recent work has shown that neglecting practical identifiability can
severely hamper epidemic predictions~\citep{gallo2022}.

We formulate the epidemic model in discrete time, matching the natural
resolution of weekly surveillance data. This yields exact, closed-form
sensitivity equations that enable efficient computation of the Fisher
Information Matrix and CRLB in a single forward pass alongside the
state equations. We apply this framework to a two-region SIR model
with seasonal forcing, deriving the complete sensitivity equations for
the connectivity parameter $\theta$.

Our central finding is that \textbf{identifiability of $\theta$
  depends on the degree of symmetry between the two regions}. When
both regions experience epidemics in phase \emph{and} start from
similar initial conditions, their incidence trajectories are nearly
indistinguishable regardless of the value of $\theta$, and the CRLB
diverges. Any source of asymmetry---different seasonal phasing or
different initial susceptibility---breaks this degeneracy and restores
identifiability. The practical danger is that a maximum-likelihood
optimizer will still converge to some $\hat\theta$ in the
symmetric case; only the CRLB reveals that this estimate is
meaningless.

We validate these predictions on synthetic data and then apply the
framework to pneumonia and influenza (P\&I) excess mortality across
the 20 largest US states.  Despite 12 seasons of data, connectivity
estimates for nearly all state pairs cannot be distinguished from
zero. This is consistent with the strong epidemic synchronization
driven by shared winter forcing across temperate North America:
seasonal phases are nearly identical, and population-level
susceptibility follows similar patterns from state to state, placing
the data squarely in the regime where the CRLB is largest.
